\chapter{Conclusion}\label{ch:conclusion}

This dissertation examined how a knowledge-management system can use machine learning to support cybersecurity incident-response decisions. We developed Sentinel to connect maintained articles and response playbooks with contextual retrieval and source-based explanations. The starting point was a practical knowledge-management problem: an analyst needs the conditions and rationale behind a procedure, not only the document containing it. We therefore treated capture, maintenance and reuse as parts of the same task, with the analyst retaining responsibility for action. Optimisation means reducing the effort needed to reach a correct, well-supported decision, rather than simply producing an answer more quickly.

\section{Answers to the research questions}
For RQ1, our design answer is to retain procedures, threat-related material and incident lessons with the context needed to judge their applicability. A playbook needs conditions, responsibilities and verification steps; an incident lesson needs the observations and reasoning behind an earlier decision. In Sentinel, classified articles, source attribution, revisions, publication review and expiry provide a way to maintain that knowledge. Retained uploads and recoverable processing extend acquisition without making every submitted document approved organisational guidance. Saved conversations likewise remain separate from published knowledge, so reuse includes an editorial decision rather than automatic promotion of generated text.

RQ2 connects this knowledge to the analyst's question. We combine embeddings from pretrained models with lexical retrieval and indicator lookup, allowing semantic similarity and explicit technical terms to contribute to passage selection. The language model receives authorised source material and instructions to distinguish reported observations from possible explanations, relate actions to their purpose and preserve uncertainty. Article and passage references connect the response to inspectable knowledge. Machine learning supports representation and explanation; it does not determine publication authority or execute remediation. When generation is unavailable, eligible playbook excerpts or direct article references provide an alternative route to the existing guidance. This design uses pretrained models at inference time rather than training a new incident-response model.

For RQ3, the documented engineering checks establish selected processing, access and interface behaviours under their test conditions. The live provider calls reported in Chapter~\ref{ch:evaluation} distinguish accepted requests from an embedding configuration rejected at the requested dimension. These observations concern technical operation; they do not quantify retrieval relevance, answer correctness or analyst efficiency. We distinguish those outcomes from usability and interpret task time alongside the quality of the resulting decision. A resolved citation, a completed model request and a correct incident-response decision answer different evaluation questions.
% RQ3 quantitative synthesis: integrate the transferred records of the completed controlled model comparisons and human evaluation here, consistently with Chapter 4. Retain actual conditions, case and participant denominators, failures, scoring rules and uncertainty; report only outcomes present in those records.
% Reserve the effectiveness answer and any comparative claim about retrieval, response quality, task efficiency or usability until that synthesis is available. Missing transferred records do not mean the evaluations were not conducted. Do not infer model rankings, participant findings or effect sizes from technical checks, compatibility calls or illustrative examples.

\section{Contributions and limits}
The principal contribution is the connection between maintained knowledge and an incident question. An analyst can describe suspicious activity, consult retrieved guidance, inspect its supporting passage and return to the broader procedure. Knowledge custodians can revise that procedure without retraining the language model. By bringing these activities together, we give knowledge management a direct role in response assistance, while keeping responsibility for the underlying content visible. The contribution is an integration of established methods around incident-response work, rather than a new retrieval algorithm.

The design also makes the distinction between source access and advice review concrete. The reading interface exposes both the article and the particular passage associated with an answer. This supports checking a recommendation against its context, but an exact reference does not prove that every generated statement follows from the passage. Nor does a relevant source establish that its procedure applies to the affected environment. The analyst still needs to assess prerequisites, operational consequences and uncertainty before acting.

Several choices constrain the context available for that assessment. Fixed passage boundaries and a limited generation context can separate an instruction from a qualification. Literal-question retrieval can miss a follow-up that needs conversational context. Publication and indexing occur separately, so an approved revision may temporarily be unavailable to semantic search. Current access rules can also make a previously cited passage unavailable. We accept these last two restrictions to avoid reusing obsolete or no longer authorised material. The graph assertion store records evidence-linked relationships; the assistant does not traverse it to retrieve knowledge.

Deployment introduces a separate set of boundaries. Local chat or embedding inference does not establish that storage and identity services are local, and an external model service receives the material sent to it. Access within the KMS and permission to disclose knowledge to a provider are therefore distinct decisions. Operational records also remain best-effort rather than a guarantee of complete event capture. These limits affect how the system can be used and investigated; they are not measures of model quality or substitutes for assessing decisions in their operational setting.

\section{Further development}
We see useful extensions in preserving procedural context during passage construction and improving follow-up retrieval without bypassing evidence eligibility. Structure-aware passages could keep prerequisites with the actions they qualify. A revised conversational search strategy could use earlier observations while retaining a clear response when accessible evidence is insufficient. Both extensions address specific ways in which relevant knowledge can be lost between a source and an answer.

A more direct route from a reviewed incident lesson to a proposed article revision could also make experience easier to retain, while leaving publication with the responsible custodians. The aim would be to preserve why a decision was made and when it applies, not to treat a saved conversation as automatically learned organisational knowledge. Linking an exception to the procedure it changes would help the next responder identify the relevant qualification.

Sentinel provides an implemented approach to the architectural part of the principal research question: maintained knowledge supplies the content, machine learning supports its retrieval and explanation, and analysts decide how to use it. The criterion for optimisation remains the quality of that decision together with the effort required to reach it. This keeps incident-response knowledge reuse, rather than the fluency of generated text, at the centre of the contribution.
% Reconcile the closing synthesis with the transferred RQ3 results after their integration, retaining this distinction between implemented capability and measured effectiveness.
